% Part: first-order-logic
% Chapter: introduction
% Section: formulas

\documentclass[../../../include/open-logic-section]{subfiles}

\begin{document}

\olfileid{fol}{int}{fml}

\section{\usetoken{P}{formula}}

د لومړۍ درجې منطق د !!{sentence}s د سخت تعريف او د !!{variable}s
له کارولو څخه راولاړو شوو مسئلو د هوارولو دپاره به لاندې لاره
کاروو۔ لومړے د عبارتونو يو \emph{بل} سټ تعريفوو: !!{formula}s۔ له
دې وروسته پکښې د !!{variable}s ونډه څېړلے شو---او د ځينو نورو
مفکورو، يعنې ``ازادو'' او ``تړلو'' !!{variable}s، پر بنسټ تعريفولے
شو چې !!a{sentence} څه ده (يعنې داسې !!a{formula} چې ازاد
!!{variable}s نۀ لري)۔ دا کار يوازې د ``!!{sentence}'' د تعريف د
اسانه کولو دپاره نۀ کوو؛ بلکې د صدق د معنايي مفهوم د تعريف په
طريقه کښې به هم ډېر اړين وي۔

راځئ ``!!{formula}'' د لومړۍ درجې د يوې ساده ژبې دپاره تعريف کړو:
داسې ژبه چې يوازې يو !!{predicate}~$\Obj P$، يو
!!{constant}~$\Obj a$، او يوازې منطقي سمبولونه $\lnot$، $\land$
او~$\lexists$ لري۔ زموږ بشپړ تعريفونه به ډېر عام وي: نامتناهي ډېر
!!{predicate}s او !!{constant}s به جايز کړو۔ په حقيقت کښې هغه
!!{function}s به هم څېړو چې له !!{constant}s او !!{variable}s سره
يوځاے ``ترمونه'' جوړوي۔ اوس دپاره~$\Obj a$ او متغيرونه زموږ يوازيني
ترمونه دي۔ البته نامتناهي ډېرو !!{variable}s ته اړتيا لرو۔ په رسمي
ډول به $\Obj v_0$، $\Obj v_1$، \dots، د متغيرونو په توګه کاروو۔

\begin{defn}
د \emph{!!{formula}s} سټ~$\Frm$ داسې تعريفېږي:
\begin{enumerate}
\item\ollabel{fmls-atom} $\Atom{\Obj P}{\Obj a}$ او $\Atom{\Obj
  P}{\Obj v_i}$ !!{formula}s دي ($i \in \Nat$)۔

\tagitem{prvNot}{\ollabel{fmls-not}که $!A$ !!a{formula} وي، نو $\lnot !A$
  !!{formula} ده۔}{}

\tagitem{prvAnd}{که $!A$ او $!B$ !!{formula}s وي، نو $(!A \land
  !B)$ !!a{formula} ده۔}{}

\tagitem{prvEx}{\ollabel{fmls-ex}که $!A$ !!a{formula} او $x$
  !!a{variable} وي، نو $\lexists[x][!A]$ !!a{formula} ده۔}{}

\tagitem{limitClause}{\ollabel{fmls-limit}بل هېڅ څيز !!a{formula} نۀ دے۔}{}
\end{enumerate}
\end{defn}

\olref{fmls-atom} راته وايي چې د هر $i \in \Nat$ دپاره
$\Atom{\Obj P}{\Obj a}$ او $\Atom{\Obj P}{\Obj v_i}$
!!{formula}s دي۔ دا تش په نامه \emph{اټومي} !!{formula}s دي۔ هغوئ
د پيل دپاره څه راکوي۔ نور شرطونه له هغو څخه چې لا مو
جوړ کړي، د نويو !!{formula}s د جوړولو لارې راکوي۔ نو مثلاً د
\olref{fmls-not} له مخې $\lnot \Atom{\Obj P}{\Obj v_2}$
!!a{formula} ده، ځکه $\Atom{\Obj P}{\Obj v_2}$ د
\olref{fmls-atom} له مخې لا !!a{formula} ده۔ بيا د
\olref{fmls-ex} له مخې $\lexists[\Obj v_2][\lnot \Atom{\Obj P}{\Obj v_2}]$
بله !!{formula} ده، او همداسې مخکښې۔ \olref{fmls-limit} راته وايي
چې \emph{يوازې} هغه تارونه !!{formula}s ګڼل کېږي چې په همدې طريقه
ئې جوړولے شو۔ په ځانګړي ډول
$\lexists[\Obj v_0][\Atom{\Obj P}{\Obj a}]$ او
$\lexists[\Obj v_0][\lexists[\Obj v_0][\Atom{\Obj P}{\Obj a}]]$
په رښتيا !!{formula}s ګڼل کېږي، خو $(\lnot \Atom{\Obj P}{\Obj a})$
نۀ ګڼل کېږي، ځکه دباندې قوسونه ئې زيات دي۔

د !!{formula}s د تعريف دا طريقه \emph{استقرايي تعريف} بلل کېږي،
او اجازه راکوي چې د استقرايي ثبوت په هغه بڼه د !!{formula}s په باب
څه ثابت کړو چې \emph{جوړښتي استقرا} بلل کېږي۔ دا دواړه په عمومي
ډول په \olref[mth][ind][idf]{sec} او \olref[mth][ind][sti]{sec}
کښې څېړل شوي دي؛ وروسته ثبوتونو ته له ننوتلو مخکښې ئې بيا وګورئ۔
بنسټيزه مفکوره دا ده: که غواړئ ثابته کړئ چې يو څۀ د ټولو
!!{formula}s دپاره رښتيا دي، لومړے ښيئ چې د اټومي !!{formula}s
دپاره رښتيا دي؛ بيا ښيئ چې \emph{که} د هر !!{formula}~$!A$
(او~$!B$) دپاره رښتيا وي، د $\lnot !A$، $(!A \land !B)$ او
$\lexists[x][!A]$ دپاره \emph{هم} رښتيا دي۔ مثلاً په دې سره ثابتېږي
چې دا د $\lexists[\Obj v_2][\lnot \Atom{\Obj P}{\Obj v_2}]$ دپاره
رښتيا دي: له لومړۍ برخې پوهېږئ چې د اټومي
!!{formula}~$\Atom{\Obj P}{\Obj v_2}$ دپاره رښتيا دي۔ بيا له دويمې
برخې ترلاسه کوئ چې د $\lnot \Atom{\Obj P}{\Obj v_2}$ دپاره رښتيا
دي، او ورپسې بيا د $\lexists[\Obj v_2][\lnot \Atom{\Obj P}{\Obj v_2}]$
په خپله دپاره رښتيا دي۔ ځکه چې ټول !!{formula}s په استقرايي ډول له
اټومي !!{formula}s څخه جوړېږي، دا طريقه د هر فارمول دپاره کار کوي۔

\end{document}
